\chapter{て形① — Forma て (I): Secuencias}
\unidadheader{13}{て形①}{Forma て (I)}{Conectar acciones y describir secuencias}

\begin{tcolorbox}[vocabbox, title={📌 Vocabulario Nuevo}]
\begin{tabularx}{\linewidth}{llX}
\toprule
\textbf{Japonés} & \textbf{Español} & \textbf{Tipo} \\
\midrule
\vocabitem{入る}{はいる}{entrar}{v. gr.1}
\vocabitem{出る}{でる}{salir}{v. gr.2}
\vocabitem{持つ}{もつ}{sostener / tener (en mano)}{v. gr.1}
\vocabitem{貸す}{かす}{prestar (dar)}{v. gr.1}
\vocabitem{借りる}{かりる}{pedir prestado / tomar}{v. gr.2}
\vocabitem{洗う}{あらう}{lavar}{v. gr.1}
\vocabitem{開ける}{あける}{abrir (algo)}{v. gr.2}
\vocabitem{閉める}{しめる}{cerrar (algo)}{v. gr.2}
\vocabitem{起こす}{おこす}{despertar a alguien}{v. gr.1}
\vocabitem{直す}{なおす}{reparar / corregir}{v. gr.1}
\bottomrule
\end{tabularx}
\end{tcolorbox}

\begin{tcolorbox}[phrasebox, title={⚙️ Reglas de conversión a て形}]
\begin{tabular}{lll}
\toprule
\textbf{Terminación} & \textbf{→ て形} & \textbf{Ejemplo} \\
\midrule
く → & いて & \ruby{書}{か}く → \ruby{書}{か}いて \\
ぐ → & いで & \ruby{泳}{およ}ぐ → \ruby{泳}{およ}いで \\
す → & して & \ruby{話}{はな}す → \ruby{話}{はな}して \\
む/ぶ/ぬ → & んで & \ruby{飲}{の}む → \ruby{飲}{の}んで \\
る/う/つ → & って & \ruby{買}{か}う → \ruby{買}{か}って \\
gr.2 (〜る) → & て & \ruby{食}{た}べる → \ruby{食}{た}べて \\
する → & して & する → して \\
くる → & きて & くる → きて \\
\bottomrule
\end{tabular}
\end{tcolorbox}

\subsection*{🗣️ Frases Clave}

\frase{\ruby{手}{て}を\ruby{洗}{あら}ってから\ruby{食}{た}べてください。}
  {Lávate las manos y luego come.}
\frase{\ruby{宿題}{しゅくだい}をしてから\ruby{遊}{あそ}びます。}
  {Hago la tarea y después juego.}
\frase{\ruby{音楽}{おんがく}を\ruby{聞}{き}きながら\ruby{勉強}{べんきょう}します。}
  {Estudio mientras escucho música.}

\subsection*{⚙️ Gramática}

\patron{Secuencia de acciones: てから}
  {V-て + から + V2 (primero ~ y luego ~)}
  {\ej{\ruby{シャワー}{しゃわー}を\ruby{浴}{あ}びてから\ruby{寝}{ね}ます。}
    {Me ducho y luego me acuesto.}
  \ej{\ruby{日本}{にほん}に\ruby{来}{き}てから\ruby{日本語}{にほんご}を\ruby{勉強}{べんきょう}しました。}
    {Después de venir a Japón, estudié japonés.}}

\patron{Acciones simultáneas: ながら}
  {V(ます-stem) + ながら + V2 (mientras hace A, hace B)}
  {\ej{コーヒーを\ruby{飲}{の}みながら\ruby{新聞}{しんぶん}を\ruby{読}{よ}みます。}
    {Leo el periódico mientras tomo café.}}

\begin{tcolorbox}[ejerciciobox, title={📝 Ejercicios Rápidos}]
\begin{enumerate}
  \item Convierte a て形: \ruby{泳}{およ}ぐ / \ruby{持}{も}つ / \ruby{出}{で}る
  \item Une con てから: \ruby{宿題}{しゅくだい}をする → \ruby{ゲーム}{げーむ}をする
  \item Describe algo que haces simultáneamente.
\end{enumerate}
\vspace{0.4em}
\textbf{Respuestas:}
1) \ruby{泳}{およ}いで / \ruby{持}{も}って / \ruby{出}{で}て\quad
2) \ruby{宿題}{しゅくだい}をしてから\ruby{ゲーム}{げーむ}をします。\quad
3) (personal)
\end{tcolorbox}

\begin{tcolorbox}[kanjibox, title={🀄 Kanjis de la Unidad}]
\begin{tabularx}{\linewidth}{cllXX}
\toprule
\textbf{Kanji} & \textbf{On} & \textbf{Kun} & \textbf{Significado} & \textbf{Vocab} \\
\midrule
\kanjirow{入}{にゅう}{はい・い}{entrar}{\ruby{入}{はい}る / \ruby{入口}{いりぐち}}
\kanjirow{出}{しゅつ}{で・だ}{salir / dar}{\ruby{出}{で}る / \ruby{出口}{でぐち}}
\kanjirow{持}{じ}{も}{sostener / tener}{\ruby{持}{も}つ / \ruby{持参}{じさん}}
\kanjirow{借}{しゃく}{か}{pedir prestado}{\ruby{借}{か}りる / \ruby{借金}{しゃっきん}}
\kanjirow{洗}{せん}{あら}{lavar}{\ruby{洗}{あら}う / \ruby{洗濯}{せんたく}}
\bottomrule
\end{tabularx}
\end{tcolorbox}
